\section{答案与解析：结构力学基础与技巧班第2次作业答案（理论力学回顾）}
\begin{analysisbox}
支座反力题统一从整体隔离体入手，列水平力、竖向力和力矩平衡。若求得反力为负，说明真实方向与假设方向相反，数值仍可直接代入后续内力计算。
\end{analysisbox}
\subsection*{第 1 页转写}
（a）解：取整体隔离体竖向力平衡，可得B处竖向支反力



0

0

1

1

2

2

By

F

l

q

q l

=

\textbackslash{}times \textbackslash{}times

=



对B点取矩得B处支座反力矩

2

0

0

1

1

1

2

3

6

B

M

l

q

l

q l

=

\textbackslash{}times \textbackslash{}times

\textbackslash{}times

=

（顺时针）

（b）解：取整体隔离体竖向力平衡，可得A处竖向支反力



15 1

30

45

Ay

F

KN

=

\textbackslash{}times +

=



对A点取矩得A处支座反力矩

15 1 2.5

30 3

127.5

A

M

KN m

=

\textbackslash{}times \textbackslash{}times

+

\textbackslash{}times =



（逆时针）

（c）解：取整体隔离体对C点取矩，由

3 30 40 1.5

0

B

F \textbackslash{}times +

-

\textbackslash{}times

=

，可得B点竖向支

反力

10

B

F

KN

=

，竖向力平衡得

30

C

F

KN

=

。

（d）解：取整体隔离体对A点取矩，由

10 4 0.2 10 5 0

B

F \textbackslash{}times

--

\textbackslash{}times

\textbackslash{}times =，可得B点竖向

支反力

1.4

B

F

KN

=

，竖向力平衡得

0.6

A

F

KN

=

。

（e）解：取整体隔离体对A点取矩，由

3 6 20 2 20

0

B

F \textbackslash{}times --

\textbackslash{}times -

=

，可得B点竖向

支反力

22

B

F

KN

=

，竖向力平衡得

2

A

F

KN

=

。
\subsection*{第 2 页转写}
（f）解：取整体隔离体对A点取矩，由

4

4.5

0

C

F

a q a

a

\textbackslash{}times

-\textbackslash{}times \textbackslash{}times

=

，可得B点竖向支

反力

1.125

C

F

qa

=

，竖向力平衡得

0.125

A

F

qa

=

。

（g）解：取整体隔离体竖向力平衡，可得A处竖向支反力



5 1

5

Ay

F

KN

=

\textbackslash{}times =



对A点取矩得A处支座反力矩

5 1 2

10

A

M

KN m

=\textbackslash{}times \textbackslash{}times

=



（逆时针）

（h）解：取整体隔离体竖向力平衡，可得A处竖向支反力



15

Ay

F

KN

=



对A点取矩得A处支座反力矩

10 15 1

25

A

M

KN m

=

+

\textbackslash{}times =



（逆时针）

（i）解：取整体隔离体对A点取矩，由

5

3

2.5

0

B

F

a q

a

a

\textbackslash{}times

-\textbackslash{}times

\textbackslash{}times

=，可得B点竖向支

反力

1.5

B

F

qa

=

，竖向力平衡得

1.5

A

F

qa

=

。
\subsection*{第 3 页转写}
（j）解：取整体隔离体对A点取矩，由

3

2

0

B

e

e

F

a M

M

\textbackslash{}times

+

-

=

，可得B点竖向支反

力

3

e

B

M

F

a

=

，竖向力平衡得

3

e

A

M

F

a

=

。

（k）解：取整体隔离体对A点取矩，可得

2

2

By

Bx

F

a

F

a

q

a a

\textbackslash{}times

+

\textbackslash{}times =\textbackslash{}times

\textbackslash{}times

；取BC部

分隔离体对C点取矩可得

By

Bx

F

a

F

a

\textbackslash{}times =

\textbackslash{}times ，联立解得

2

3

By

F

qa

=

，

2

3

Bx

F

qa

=

。再由

整体竖向力和水平力平衡，可得A处支座反力

2

3

Ay

F

qa

=

，

4

3

Ax

F

qa

=

。

（l）解：取整体隔离体对A点取矩，可得

10

2 5 7.5

7.5

By

By

F

F

KN

\textbackslash{}times

=\textbackslash{}times \textbackslash{}times



=

；取

BC部分隔离体对C点取矩可得

25

5

6

2 5 2.5

12

By

Bx

Bx

F

F

F

KN

\textbackslash{}times

=

\textbackslash{}times +\textbackslash{}times \textbackslash{}times



=

。再由整

体竖向力和水平力平衡，可得A处支座反力

2.5

Ay

F

KN

=

，

25

12

Ax

F

KN

=

。
\subsection*{第 4 页转写}
（m）解：取整体隔离体对A点取矩，可得

2

1

2

2

Dy

Dy

qh

F

l

q

h

h

F

l

\textbackslash{}times =

\textbackslash{}times

\textbackslash{}times



=

，再由

整体竖向力和水平力平衡，可得A处支座反力

2

2

Ay

qh

F

l

=

，

Dx

F

qh

=

。
